%قبل از شروع مسئله ترجیحاً چیزی ننویسید. اما اگر خواستید بنویسید، در قالب لیستی (فقرات و شمارش) بنویسید.
%برای اضافه کردن نمره، بعدی از شروع‌مسئله، [x] را اضافه کنید و به جای x نمره مورد نظر را قرار دهید. برای زیرمسئله‌ها هم به همین صورت است.
%\شروع‌مسئله[10]
\شروع‌مسئله [15]
میزان «شنگدباویت» تعدادی مجموعه از اعداد طبیعی 
$A_1, \ A_2, \ ..., \ A_n$
که همه‌ اعدادشان کمتر مساوی $M$ هستند برابر تعداد اعدادی از $1$ تا $M$ هست که می‌توان به صورت جمع اعضایی از $A_1$ تا $A_n$ نوشت به طوریکه دو عضو از یک مجموعه باهم انتخاب نشوند.
به طور دقیق‌تر فرض کنید $X$ را می‌توان به صورت $\sum_{i=1}^{n}  v_i$ نوشت به طوری‌که $v_i \in A_i$ میزان شنگدباویت این مجموعه‌ها برابر تعداد $X$ های طبیعی کمتر از $M$ای هست که می‌توان به صورت مجموع بیان شده نوشت.

الف) فرض کنید $n = 2$ است و دو مجموعه $A$ و $B$ را داریم. الگوریتمی بیابید که با زمان بهتر از $O(M^2)$ شنگدباویت این دو مجموعه را بیابد. [۷ نمره]

ب) حالا برای حالتی که $A_i$ ها همگی برابرند الگوریتمی با زمان بهتر از $O(M^2\times log_2n)$ برای یافتن شنگدباویت $n$ مجموعه بیابید. [۸ نمره]

راهنمایی: سعی کنید هر مجموعه را به‌صورت عددی در مبنای $M + 1$ ببینید که رقم $i$ام آن یک است اگر $i$ در مجموعه ظاهر شده‌باشد.
\\
\\

%\متن‌سیاه{سوال عملی اول:}\\
%\متن‌سیاه{دراپ‌تست}\\

%\متن‌سیاه{ورودی:}\\
%در خط اول ورودی، ابتدا k و سپس m داده می‌شود.\\
%\متن‌سیاه{خروجی:}\\
%در تنها سطر خروجی، تعداد آیفون‌های مورد نیاز می‌آید. اگر در هیچ صورتی چنین کاری ممکن نبود ۱- چاپ می‌شود.\\
%\متن‌سیاه{محدودیت‌ها:}\\
%$1\leq k\leq 10^4$

%$m\leq10^9$
%\\
%\textbf{\متن‌سیاه{ورودی و خروجی نمونه}\\
%% تست اول
%\begin{latin}
% \begin{tabularx}{\textwidth}{|X | X|}
% \hline
% \textbf{stdin} &
% \textbf{stdout}\\
% \hline
% 100 100&
% 1\\
% \hline
% 
% \end{tabularx}
%\end{latin}
%
%% تست دوم
%
%\begin{latin}
% \begin{tabularx}{\textwidth}{|X | X|}
% \hline
% \textbf{stdin} &
% \textbf{stdout}\\
% \hline
% 340 20&
% 3\\
% \hline
% 
% \end{tabularx}
%\end{latin}}
% صورت اصلی مسئله را در این‌جا بنویسید.



\solution{
	الف) اگر هر مجموعه $A$ و $B$  را به صورت دو عدد به صورتی ببینیم که در مبنای $M+1$ هستند و رقم $i$ ام آن یک است اگر آن عضو در مجموعه آمده باشد می‌توانیم دو عدد را در هم ضرب کنیم. اگر $C$ را عدد حاصل ضرب این‌ دو بگیریم و به ازای هر عدد مثل $X$ رقم $i$ام آن‌را به صورت $X_i$ نشان دهیم. رابطه زیر برقرار است:
	$C_i = \sum_{j = 0}{M - i} A_j \times B_{M - j} < M + 1$
	توجه کنید که چون همه ارقام گمتر مساوی ۱ هستند و رقم سمت چپ همیشه صفر است چنین رابطه‌ای برقرار است.
	
	اکنون بیایید و هر رقم بزرگتر از ۱ در $C$ را به جایش ۱ بگذارید. اکنون اگر رقمی در $C$ برابر ۱ بود یعنی به صورت جمع دو عضو در $A$ و $B$ نوشته می‌شود و اگر برابر صفر باشد یعنی نمی‌توان آن‌را به اینصورت نوشت. اگر برای ضرب از الگوریتم کاراتسوبا یا هر ضربی بهتر از $O(M^2)$ استفاده کنیم به نتیجه می‌رسیم.
	
	ب) توجه کنید در الگوریتم قبلی توانستیم دو مجموعه را ادغام کنیم تا به مجموعه متناظر $C$ برسیم. فرض کنید هر سری که ضرب می‌کنیم در نهایت همه ارقام را حداکثر ۱ می‌کنیم و با این تعریف ضرب مقدار مجموعه را به توان $n$ می‌رسانیم. صرفا کافیست از نکته زیر استفاده کنیم تا الگوریتمی بازگشتی با مرتبه زمانی بهتر از $O(M^2 \times log_2n)	$ بیابیم:
	
	$$A^{2n} = (A^n)^2 \ \to 2 * o(M^2)$$
	$$A^{2n+1} = (A^n)^2 \times A \ \to 3 * o(M^2)$$
	
}

%%%%%%%%%%%%%%%%%%%%%%%%%%%%%%%%%%%%%%%%%%%%%%%%%%%%%
% برخی دستورات مفید برای تمرین‌ها. %
%%%%%%%%%%%%%%%%%%%%%%%%%%%%%%%%%%%%%%%%%%%%%%%%%%%%%

% فرمول‌های نمونه:
%$f(x) = \Theta(x)$
%$T(x, y) \leq \log(x) + 2^y + \binom{2n}{n}$
% {\tt code\_function}


% لیست موارد بدون شماره‌گذاری:
%\شروع{فقرات}
% \فقره مورد اول 
% \شروع{فقرات}
% \فقره مورد اول داخلی 
% \شروع{فقرات}
% \فقره اولین مورد خیلی داخلی 
% \فقره دومین مورد خیلی داخلی 
% \پایان{فقرات}
% \فقره مورد دوم داخلی 
% \پایان{فقرات}
% \فقره مورد دوم 
%\پایان{فقرات}

% لیست موارد شماره‌گذاری شده:
%\شروع{شمارش}
% \فقره یک مورد 
% \فقره یک مورد دیگر 
% \فقره باز هم یک مورد دیگر
%\پایان{شمارش}

% برای نوشتن یک فرمول چند خطی که یک شماره بخوره. برای حالت‌های دیگر به سایت زیر مراجعه کنید:
% https://www.overleaf.com/learn/latex/Aligning_equations_with_amsmath
%\begin{align} 
% \begin{split}
% (x+y)^3 &= (x+y)^2(x+y) \\
% &= (x^2+2xy+y^2)(x+y) \\
% &= (x^3+2x^2y+xy^2) + (x^2y+2xy^2+y^3) \\
% &= x^3+3x^2y+3xy^2+y^3
% \end{split} 
%\end{align}

% نوشتن ماتریس و آرایه دوبعدی
%\begin{align}
% A = 
% \begin{bmatrix}
% A_{11} & A_{21} \\
% A_{21} & A_{22}
% \end{bmatrix}
%\end{align}

% اضافه کردن تصویر با شرح. برای اضافه کردن تصویر بدون شرح، فقط دستور درج‌تصویر را استفاده کنید.
%\شروع{شکل}[h!]
%\تنظیم‌ازوسط
% \درج‌تصویر[پهنا=0.5\پهنای‌متن]{img.jpg}
%\شرح{شرح تصویر}
%\برچسب{تصویر:۱}
%\پایان{شکل}
%\پایان{زیرمسائل}

%کد زیر را ببینید.
%%\lstinputlisting{source_filename.py}
%\begin{latin}
%\begin{lstlisting}
%# Import the modules
%import sys
%import random
%
%ans = True
%
%while ans:
%question = raw_input("Ask the magic 8 ball a question: (press enter to quit) ")
%
%answers = random.randint(1,8)
%
%if question == "":
%sys.exit()
%
%elif answers == 1:
%print "It is certain"
%
%elif answers == 2:
%print "Outlook good"
%
%elif answers == 3:
%print "You may rely on it"
%
%elif answers == 4:
%print "Ask again later"
%
%elif answers == 5:
%print "Concentrate and ask again"
%
%elif answers == 6:
%print "Reply hazy, try again"
%
%elif answers == 7:
%print "My reply is no"
%
%elif answers == 8:
%print "My sources say no"
%\end{lstlisting}
%\end{latin}